\section{量化理论的公理}

\begin{definition}{量化理论(一阶理论)}{}
	包含如下公理模式导出的所有隐式公理的逻辑理论$\mathcal{T}$称为一个量化理论.
	\begin{center}
		设$\boldsymbol{R}$是理论$\mathcal{T}$中的公式,$\boldsymbol{T}$是$\mathcal{T}$中的项,$\boldsymbol{x}$是一个字母,则$\left(\boldsymbol{T}|\boldsymbol{x}\right)\boldsymbol{R}\implies\left(\exists\boldsymbol{x}\right)\boldsymbol{R}$是一个公理.
	\end{center}
	该公理模式称为存在实例化公理模式.
\end{definition}

\begin{remark}{}{}
	上述公理模式表明,如果存在一个对象$\boldsymbol{T}$,使得被视作表达$\boldsymbol{x}$之性质的表达式$\boldsymbol{R}$成立,那么$\boldsymbol{R}$对于对象$\tau_{\boldsymbol{x}}\left(\boldsymbol{R}\right)$也成立,这与我们在前文中赋予$\tau_{\boldsymbol{x}}\left(\boldsymbol{R}\right)$的直观含义是吻合的.
\end{remark}

\begin{theorem}{存在实例化公理的代入封闭性}{}
	存在实例化公理模式是合法的公理模式.
\end{theorem}